\chapter{Logistic Regression and Multiclass Classification}
\companion{StatQuest: Logistic Regression}{\ytid{yIYKR4sgzI8}}

Logistic regression was the model the InfoEdge Round 3 candidate trained in the coding round, and it is
still the baseline for click, apply and lead-conversion prediction across the industry. Know it well enough
to derive its gradient on paper and to interpret every coefficient.

\section{Why not just use a straight line for yes/no?}

Suppose $y = 1$ if a seeker applies to a job we show them and $0$ otherwise. We would like a model that outputs
the \emph{probability} of applying. A straight line $w^\top x + b$ can output $-3$ or $7$, which are not
probabilities. We need a function that takes any real number and squeezes it into $(0,1)$.

\section{The sigmoid: squeezing a score into a probability}

\[ \sigma(z) = \frac{1}{1 + e^{-z}}. \]
\begin{itemize}
\item $z$ very negative: $e^{-z}$ is huge, so $\sigma(z)\approx 0$.
\item $z = 0$: $\sigma(0) = 1/2$.
\item $z$ very positive: $e^{-z}\approx0$, so $\sigma(z)\approx 1$.
\item Handy values: $\sigma(1) = 0.731$, $\sigma(2) = 0.881$, $\sigma(-2) = 0.119$, $\sigma(\ln 3) = 0.75$.
\item Its derivative is beautifully simple: $\sigma'(z) = \sigma(z)(1-\sigma(z))$, largest (0.25) at $z = 0$.
\end{itemize}

\begin{center}
\begin{tikzpicture}
\begin{axis}[width=0.6\linewidth,height=4.2cm,xlabel={$z = w^\top x + b$},ylabel={$\sigma(z)$},
  xmin=-6,xmax=6,ymin=0,ymax=1,ytick={0,0.5,1},grid=major,grid style={black!8}]
\addplot[domain=-6:6,samples=80,very thick,accent]{1/(1+exp(-x))};
\end{axis}
\end{tikzpicture}
\end{center}

\textbf{Logistic regression} computes a linear score $z = \mathbf w^\top\mathbf x + b$ and outputs
$p = \sigma(z) = P(y = 1\mid \mathbf x)$.

\section{Odds and log-odds: how to read the coefficients}

The \textbf{odds} of an event with probability $p$ are $p/(1-p)$. Probability 0.8 means odds $4$ (``4 to 1'').
Inverting the sigmoid gives
\[ \log\frac{p}{1-p} = \mathbf w^\top\mathbf x + b. \]
The left side is the \textbf{log-odds} or \textbf{logit}. So logistic regression is a \emph{linear model for
the log-odds}. This gives the standard interpretation: \emph{increasing $x_j$ by one unit adds $w_j$ to the
log-odds, i.e. multiplies the odds by $e^{w_j}$}, holding other features fixed.

\begin{example}[: interpreting a coefficient]
In an apply-prediction model, the feature ``job is in the seeker's preferred city'' (0/1) has $w = 1.1$.
$e^{1.1} = 3.0$: the odds of applying triple. If a seeker's base probability was $0.10$ (odds $0.111$), the
new odds are $0.333$ and the new probability is $0.333/1.333 = 0.25$. The probability did not triple; the
\emph{odds} did.
\end{example}

\section{The decision boundary is a straight line}
To turn a probability into a class we use a threshold $t$, often $0.5$. Predict 1 when $\sigma(z)\ge0.5$, i.e.
when $z \ge 0$. The set $\mathbf w^\top\mathbf x + b = 0$ is a line (a hyperplane in higher dimensions). So,
despite the curved sigmoid, logistic regression is a \textbf{linear classifier}. A different threshold $t$ just
moves the line parallel to itself ($z \ge \log\frac{t}{1-t}$). Curved boundaries need engineered features
($x_1^2$, $x_1x_2$).

\section{Training: maximum likelihood and log loss}

\subsection{From probabilities to a loss}
For one example, the model assigns probability $p$ to $y = 1$ and $1-p$ to $y = 0$. The probability it assigns
to what actually happened is $p^y(1-p)^{1-y}$ (this is $p$ if $y=1$ and $1-p$ if $y = 0$). Assuming examples
are independent, the probability of the whole dataset is the product of these. \textbf{Maximum likelihood} picks
the weights that make the observed data most probable. Products are awkward, so take logs (which turns
products into sums and does not move the maximum) and flip the sign to get something to minimise:
\[ L(\mathbf w) = -\frac1n\sum_{i=1}^n\Big[y_i\log p_i + (1-y_i)\log(1-p_i)\Big]. \]
This is the \textbf{log loss}, also called \textbf{binary cross-entropy}.

\begin{example}[: how log loss punishes confidence]
True label $y=1$. If $p = 0.9$: loss $= -\ln 0.9 = 0.105$. If $p = 0.5$: $0.693$. If $p = 0.1$: $2.303$.
If $p = 0.01$: $4.605$. Being confidently wrong is punished brutally, which forces honest probabilities.
\end{example}

\subsection{The gradient, derived}
With $p = \sigma(z)$ and $z = \mathbf w^\top\mathbf x + b$, for one example
$\ell = -[y\log p + (1-y)\log(1-p)]$. By the chain rule:
\[ \frac{\partial\ell}{\partial p} = -\frac{y}{p} + \frac{1-y}{1-p},\qquad \frac{\partial p}{\partial z} = p(1-p),\qquad
\frac{\partial z}{\partial w_j} = x_j. \]
Multiplying, $\frac{\partial\ell}{\partial z} = -y(1-p) + (1-y)p = p - y$, so
\[ \boxed{\frac{\partial\ell}{\partial w_j} = (p - y)\,x_j,\qquad \frac{\partial\ell}{\partial b} = p - y.} \]
``Prediction minus truth, times the input'': exactly the same form as linear regression's gradient. The $p(1-p)$
from the sigmoid cancelled perfectly; this is why log loss and sigmoid are a natural pair.

\subsection{No closed form, but a convex problem}
Setting the gradient to zero gives equations with no algebraic solution, so we iterate: gradient descent, or
Newton's method (called \textbf{IRLS}, iteratively reweighted least squares, because each Newton step is a
weighted least-squares fit). The log loss is \textbf{convex} in $\mathbf w$: shaped like a bowl, with no false
valleys, so any of these methods reaches the global optimum.

\begin{example}[: one gradient step by hand]
One feature, $w=0.5$, $b=0$, learning rate $\eta = 0.1$, example $x = 2$, $y = 0$.
$z = 1$, $p = \sigma(1) = 0.731$. Gradients: $(p-y)x = 1.462$ and $p - y = 0.731$.
Update: $w = 0.5 - 0.1462 = 0.354$, $b = -0.073$. Both moved to lower $p$, correctly, since the true label is 0.
\end{example}

\subsection{Why not squared error with a sigmoid?}
\begin{enumerate}
\item It is non-convex in $\mathbf w$: there can be flat plateaus and local minima.
\item Its gradient contains $p(1-p)$, which is near zero when $p\approx0$ or $1$. A confidently wrong model
  ($p=0.001$ for a positive) barely learns: the gradient vanishes exactly when the error is largest.
\item Log loss is the maximum-likelihood loss for Bernoulli labels; squared error corresponds to Gaussian noise,
  which binary labels do not have.
\end{enumerate}

\section{Regularisation and separable data}
If the classes can be perfectly separated by a line, the log loss can be pushed toward zero by making $\norm{\mathbf w}$
larger and larger (the sigmoid gets steeper; every probability approaches 0 or 1). Unregularised training never
converges: weights grow without bound. An L2 penalty fixes this. scikit-learn's \code{LogisticRegression} applies L2
by default with \code{C=1.0}, where $C = 1/\lambda$: \emph{smaller $C$ means stronger regularisation}.

\section{Probabilities and calibration}
A model is \textbf{calibrated} if, among cases where it says 0.3, about 30\% are positive. Logistic regression fitted
by maximum likelihood tends to be well calibrated when its features are adequate, which is a big reason it is still used
for bidding, ranking and expected-value decisions (``expected applies $=$ sum of probabilities''). Two things break
calibration: heavy regularisation, and resampling the classes (training on 50/50 data when reality is 1/99 shifts the
intercept; correct it with $b \leftarrow b + \log\frac{\pi_{true}/(1-\pi_{true})}{\pi_{train}/(1-\pi_{train})}$).

\section{More than two classes}
\subsection{Softmax (multinomial) regression}
One score per class, $z_k = \mathbf w_k^\top\mathbf x + b_k$, turned into probabilities that sum to one:
\[ p_k = \frac{e^{z_k}}{\sum_j e^{z_j}}. \]
Loss: $-\log p_{\text{true class}}$ (categorical cross-entropy).
\begin{example}[: softmax by hand]
Scores $(2, 1, 0)$: $e^2 = 7.389$, $e^1 = 2.718$, $e^0 = 1$, sum $11.107$. Probabilities
$(0.665, 0.245, 0.090)$. If the true class is the second, loss $= -\ln 0.245 = 1.41$.
Subtracting the maximum score before exponentiating gives the same probabilities and avoids overflow.
\end{example}
With two classes softmax reduces exactly to the sigmoid on $z_1 - z_2$.

\subsection{One-vs-Rest and One-vs-One}
For any binary classifier:
\begin{itemize}
\item \textbf{One-vs-Rest (OvR, OvA)}: $K$ classifiers, ``class $k$ vs everything else''; predict the class with the
  highest score. Each model sees all data but is imbalanced.
\item \textbf{One-vs-One (OvO)}: $\binom{K}{2} = K(K-1)/2$ classifiers, one per pair, each trained only on those two
  classes; predict by majority vote. Used by kernel SVMs because many small problems are cheaper than a few large ones.
\end{itemize}
For 5 classes: OvR trains 5 models, OvO trains 10.

\section{Multi-label is different}
\textbf{Multiclass}: exactly one label (a job belongs to one function: IT, Sales, \dots). \textbf{Multi-label}:
any number of labels (a resume can have many skills). Multi-label uses an independent sigmoid per label with binary
cross-entropy each, not a softmax.

\begin{practical}[: logistic regression at InfoEdge scale]
\begin{itemize}
\item \textbf{Apply / click prediction} in the Apply Generation engine and job recommendation: billions of rows,
  sparse one-hot and hashed features, logistic regression trained with SGD. Fast to score, calibrated, and the
  coefficients explain behaviour.
\item \textbf{99acres telecalling}: ``probability this user is a high-intent buyer'' ranks whom to call first.
  Calibration lets the business compute expected conversions per caller-hour.
\item \textbf{Spam / fake profile detection} as a binary classifier, where the panel explicitly said to look at recall.
\end{itemize}
\end{practical}

\stuck{StatQuest: Logistic Regression Details Pt 2, Maximum Likelihood}{\yts{StatQuest+Logistic+Regression+Details+Pt+2+Maximum+Likelihood}}
\section{Interview questions}

\begin{qa}{\pyq{Round 3 coding round} Walk me through logistic regression: what it models, how it is trained,
and how you would evaluate it on the placement dataset.}
\oneline{It models the log-odds of the positive class as a linear function of the features, is trained by minimising
log loss (maximum likelihood) with gradient-based optimisation, and is evaluated with a confusion matrix, precision,
recall, F1 and ROC AUC on a held-out set.}
\why
\textbf{Model.} Score $z = \mathbf w^\top\mathbf x + b$, probability $p = \sigma(z)$. Equivalently
$\log\frac{p}{1-p} = \mathbf w^\top\mathbf x + b$.
\textbf{Training.} Minimise $-\frac1n\sum[y\log p + (1-y)\log(1-p)]$ plus an L2 penalty. The gradient is
$(p - y)\mathbf x$; the problem is convex, so the optimiser (lbfgs by default in scikit-learn) finds the global minimum.
\textbf{Workflow on the placement data.} Drop the leaking \code{salary} column; one-hot encode board, stream and work
experience; scale the percentage columns (the penalty treats all weights alike); stratified train/test split; fit;
predict; print confusion matrix and classification report; look at ROC AUC; choose a threshold to suit the cost of
errors.
\worked
\begin{lstlisting}
X = df.drop(columns=["status", "salary"])          # salary leaks the target
y = (df["status"] == "Placed").astype(int)
num = X.select_dtypes("number").columns; cat = X.columns.difference(num)
pre = ColumnTransformer([("num", StandardScaler(), num),
                         ("cat", OneHotEncoder(handle_unknown="ignore"), cat)])
clf = Pipeline([("pre", pre), ("lr", LogisticRegression(max_iter=1000))])
Xtr, Xte, ytr, yte = train_test_split(X, y, test_size=0.2, stratify=y, random_state=42)
clf.fit(Xtr, ytr); pred = clf.predict(Xte)
print(confusion_matrix(yte, pred)); print(classification_report(yte, pred))
print(roc_auc_score(yte, clf.predict_proba(Xte)[:, 1]))
\end{lstlisting}
\pushes \emph{``How do you read the coefficients?''} After scaling, $e^{w_j}$ is the odds multiplier per one standard
deviation of feature $j$. \emph{``Why scale?''} For the penalty's fairness and for faster convergence.
\end{qa}

\begin{qa}{Derive the gradient of the log loss for logistic regression.}
\oneline{By the chain rule the derivative with respect to the score is $p - y$, so the gradient with respect to $w_j$ is
$(p - y)x_j$.}
\why Written out in the theory section: $\partial\ell/\partial p = -y/p + (1-y)/(1-p)$; $\partial p/\partial z =
p(1-p)$; product $= -y(1-p) + (1-y)p = p - y$; times $\partial z/\partial w_j = x_j$. For the whole dataset,
$\nabla_{\mathbf w}L = \frac1n X^\top(\mathbf p - \mathbf y)$.
\pushes \emph{``What is the Hessian?''} $\frac1n X^\top S X$ with $S = \text{diag}(p_i(1-p_i))$. It is positive
semi-definite, which proves convexity, and it is what Newton/IRLS uses.
\end{qa}

\begin{qa}{Why is it called logistic \emph{regression} if it is a classifier?}
\oneline{Because it regresses (fits a linear model for) the log-odds of the outcome, producing a continuous probability;
it becomes a classifier only when you add a threshold.}
\why It is a generalised linear model: a linear predictor, a link function (logit) and a Bernoulli distribution
(Chapter 4). The output is a probability, a continuous quantity; classification is a decision we layer on top.
\end{qa}

\begin{qa}{A coefficient for ``has premium account'' is $0.693$. Interpret it.}
\oneline{Premium users have twice the odds of applying ($e^{0.693} = 2$), other features held fixed; this is not a doubling
of probability.}
\why Log-odds increase by 0.693, so odds multiply by 2. From $p = 0.5$ (odds 1) the new odds are 2 and $p = 2/3$; from
$p=0.05$ (odds 0.0526) the new odds are 0.105 and $p = 0.095$. Near small probabilities, odds ratios approximate risk
ratios.
\end{qa}

\begin{qa}{Why do we not use mean squared error to train logistic regression?}
\oneline{Because with a sigmoid output MSE is non-convex and its gradient vanishes when the model is confidently wrong,
whereas log loss is convex, is the maximum-likelihood loss for binary labels, and gives the clean gradient $(p-y)x$.}
\why The MSE gradient is $(p - y)\,p(1-p)\,x$. For a positive example predicted $p = 0.001$, $p(1-p)\approx0.001$, so
the update is a thousand times weaker than it should be. Log loss cancels that factor.
\worked $y = 1$, $p = 0.01$: log-loss gradient w.r.t.\ $z$ is $-0.99$; MSE gradient is $-0.99\cdot0.0099\approx -0.0098$.
\end{qa}

\begin{qa}{What happens when the training data are perfectly separable?}
\oneline{The unregularised maximum-likelihood solution does not exist: weights grow without bound as the optimiser keeps
steepening the sigmoid; regularisation gives a finite answer.}
\why If a line separates the classes, scaling $\mathbf w$ by $c>1$ keeps the same boundary but makes every probability
more extreme and every loss term smaller, so the loss keeps decreasing toward 0. Software prints convergence warnings
and coefficients with huge standard errors. L2 penalty, early stopping, or Firth's correction fixes it.
\end{qa}

\begin{qa}{\deep Is logistic regression a linear model? Can it learn XOR?}
\oneline{It is linear in its decision boundary (a hyperplane in feature space), so on raw features it cannot learn XOR;
with an interaction feature $x_1x_2$ it can.}
\why For XOR points $(0,0)\to0,(1,1)\to0,(0,1)\to1,(1,0)\to1$, no line separates the classes. Adding $x_3 = x_1x_2$:
the score $z = 10x_1 + 10x_2 - 20x_3 - 5$ gives $-5, 5, 5, -5$ for the four points: correct. The model is still
linear in its parameters; the feature space was enriched. This is the same idea as kernels (Chapter 9) and hidden layers
(Chapter 18).
\end{qa}

\begin{qa}{\pyq{OA: one-vs-one and one-vs-all} You have 6 job functions. How many binary classifiers do OvR and OvO
train, and when would you prefer each?}
\oneline{OvR trains 6, OvO trains 15; OvR is simpler and cheaper for most models, OvO suits algorithms whose cost grows
faster than linearly with data size, like kernel SVMs.}
\why OvO: $\binom62 = 15$, each on only two classes' data (about a third of the rows here). For kernel SVMs, training
cost is between $O(n^2)$ and $O(n^3)$, so 15 small problems can be cheaper than 6 big ones. OvR gives each model a
heavily imbalanced problem (one class vs five) and its scores may not be comparable across models. For logistic
regression, use true multinomial softmax instead of either.
\end{qa}

\begin{qa}{Softmax vs independent sigmoids: when do you use which?}
\oneline{Softmax when classes are mutually exclusive (exactly one is true); independent sigmoids when labels can co-occur
(multi-label).}
\why Softmax probabilities sum to 1, so raising one lowers the others; that is right for ``which job function'' and wrong
for ``which skills does this resume contain''. For skills, use one sigmoid per skill with binary cross-entropy, and
tune a threshold per skill.
\end{qa}

\begin{qa}{How does the decision threshold affect precision and recall? How do you choose it?}
\oneline{Raising the threshold makes the model predict positive less often: precision usually rises and recall falls;
choose it from business costs, a target recall or precision, or by maximising F1 on validation data.}
\why The model's probabilities are fixed; the threshold only decides where to cut. For spam the panel wanted recall
considered; for a telecalling queue with 100 calls per day, choose the cut that fills the queue with the highest
probabilities. Formally, with cost $C_{FP}$ and $C_{FN}$ and calibrated $p$, predict positive when
$p > C_{FP}/(C_{FP} + C_{FN})$.
\worked If missing a fraudulent recruiter costs 10 times a false alarm, threshold $= 1/(1+10) = 0.09$.
\end{qa}

\begin{qa}{\deep Your training set was rebalanced from 2\% positives to 50\% by undersampling. What happens to the
predicted probabilities and how do you fix them?}
\oneline{They become inflated, because the model learned a 50\% base rate; the ranking is mostly unaffected, and you fix
the probabilities by shifting the intercept by the log of the odds ratio between true and training base rates (or by
recalibrating on untouched data).}
\why Undersampling negatives by a factor $\beta$ multiplies the odds of every prediction by $1/\beta$, which in a
logistic model only changes the intercept. With true rate $\pi = 0.02$ (odds $0.0204$) and training rate $0.5$ (odds 1),
subtract $\log(1/0.0204) = 3.89$ from $b$. Alternatively fit isotonic or Platt calibration on a validation set with the
real class balance.
\worked A raw output of $0.6$ (log-odds 0.405) corrects to log-odds $0.405 - 3.89 = -3.49$, i.e. $p = 0.030$.
\end{qa}

\begin{qa}{What is $C$ in scikit-learn's LogisticRegression? What are the defaults?}
\oneline{$C$ is the inverse regularisation strength ($C = 1/\lambda$); the default is L2 with $C = 1.0$, so smaller $C$
means stronger shrinkage.}
\why The objective is $\frac12\norm{\mathbf w}^2 + C\sum_i \ell_i$. Large $C$ trusts the data more (less regularisation),
small $C$ shrinks more. Default solver is lbfgs (L2 only); \code{liblinear} or \code{saga} support L1. Tune $C$ on a log
grid such as $10^{-3},\dots,10^{3}$ with \code{LogisticRegressionCV}.
\end{qa}

\begin{qa}{\deep Logistic regression vs a linear SVM: what is really different?}
\oneline{Both learn a linear boundary with an L2 penalty; they differ only in the loss: log loss (smooth, every point
contributes, outputs probabilities) versus hinge loss (zero for points beyond the margin, only support vectors matter,
outputs scores).}
\why Hinge $\max(0, 1 - yz)$ with $y\in\{-1,1\}$ is exactly zero once a point is correctly classified with margin, so
confidently correct points stop influencing the boundary. Log loss $\log(1+e^{-yz})$ never reaches zero, so all points
keep a small influence. In practice accuracy is often similar; logistic regression is preferred when calibrated
probabilities are needed, SVM when the margin-based robustness to far-away points helps or with kernels.
\end{qa}

\begin{qa}{What is the decision boundary of logistic regression with features $x_1$ and $x_2$ and weights $w = (2,
-1)$, $b = 1$? Classify the point $(1, 4)$.}
\oneline{The boundary is the line $2x_1 - x_2 + 1 = 0$; the point $(1,4)$ has score $-1$, probability $0.269$, so it is
class 0 at threshold 0.5.}
\why $z = 2\cdot1 - 4 + 1 = -1$; $\sigma(-1) = 1/(1 + e) = 0.269$. Points above the line $x_2 = 2x_1 + 1$ are class 0.
\end{qa}

\begin{qa}{How would you handle a categorical feature with 3{,}000 levels (for example company name) in logistic
regression?}
\oneline{One-hot encode it with regularisation (sparse is fine for linear models), or group rare levels, hash it, or use
target encoding computed out-of-fold.}
\why One-hot creates 3{,}000 sparse columns; with L2 and enough data per level this works well and is standard in click
models. Rare companies (a few rows) get weights shrunk toward 0, which is sensible. Alternatives: map to a taxonomy
(the SilverSkAI-style normalisation collapses variants like ``TCS'', ``Tata Consultancy''), an ``other'' bucket for
levels with $<20$ rows, feature hashing, or out-of-fold target encoding (never in-fold: that leaks the label).
\end{qa}

\begin{qa}{\deep Can logistic regression be used for regression of a proportion (for example the fraction of
recruiters who open a profile)?}
\oneline{Yes: fit a binomial GLM with the logit link, using counts of successes and trials, or a fractional logit;
the sigmoid keeps predictions inside $[0,1]$.}
\why The binomial GLM models $k$ successes out of $n$ trials with $P = \sigma(\mathbf w^\top\mathbf x)$; the log loss
weighted by $n$ is its likelihood. scikit-learn's \code{LogisticRegression} accepts only 0/1 labels but you can expand
rows or use \code{sample\_weight}; statsmodels' \code{GLM(family=Binomial())} accepts proportions with
\code{var\_weights}. This avoids the main weaknesses of linear regression on proportions (predictions outside $[0,1]$
and non-constant variance $p(1-p)/n$).
\end{qa}

\begin{qa}{What is the log loss of a model that always predicts the base rate $\pi$? Why is it a useful baseline?}
\oneline{It is the entropy of the labels, $-[\pi\log\pi + (1-\pi)\log(1-\pi)]$; any useful model must beat it.}
\why With $\pi = 0.1$: $-(0.1\ln0.1 + 0.9\ln0.9) = 0.230 + 0.095 = 0.325$. A model at 0.30 is only modestly
informative. Normalised cross-entropy (model log loss divided by this) is a common metric in ad/click teams.
\end{qa}

\begin{qa}{How do you get feature importance from logistic regression, and what can go wrong?}
\oneline{Use the absolute values of coefficients after standardising features (or odds ratios), and their confidence
intervals; collinearity, unscaled features and regularisation can all distort the picture.}
\why Without scaling, a coefficient depends on units. With correlated features, credit is split arbitrarily. With
L1, one of a correlated group survives. Permutation importance on validation data is a model-agnostic check.
\end{qa}

\begin{qa}{\deep Logistic regression gives probability 0.97 for a profile being fake. Should you auto-ban the user?}
\oneline{Only if the model is calibrated at that range, the cost of a false ban has been weighed, and the precision at
that threshold is verified on recent labelled data; otherwise route to human review.}
\why A probability is only meaningful if calibrated (check a reliability curve on recent data). Even at 97\% precision,
3\% of bans are wrong; for paying recruiters that may be unacceptable. A common design: auto-action above a very high
threshold with measured precision $>99.5\%$, human review between, no action below. Monitor drift, because fraud adapts.
\end{qa}

\begin{qa}{Compute the probability and log loss for $z = 1.2$ and true label 1.}
\oneline{$p = \sigma(1.2) = 0.769$, log loss $= -\ln 0.769 = 0.263$.}
\why $e^{-1.2} = 0.301$; $p = 1/1.301 = 0.769$; $-\ln(0.769) = 0.263$.
\end{qa}
